\section{Keterampilan dan Kualifikasi yang Didapatkan}

Selama masa magang di PT Inti Utama Solusindo, terdapat sejumlah keterampilan
teknis dan non-teknis yang berhasil dikembangkan secara signifikan, sejalan
dengan perubahan ruang lingkup pekerjaan dari \textit{Front-End Developer}
(Flutter) menjadi \textit{Full-Stack Developer} (React dan Toraja).

Pada periode awal (Januari--pertengahan Februari 2026), kemampuan pengembangan
UI menggunakan Flutter semakin terasah, khususnya dalam implementasi
\textit{widget} berbasis pendekatan \textit{Server-Driven UI} (SDUI),
penanganan \textit{state management}, serta \textit{debugging} pada
\textit{codebase} skala enterprise. Pengalaman melakukan migrasi
\textit{widget} dari Flutter v3.16.9 ke v3.38.5 (Task \#11858) sekaligus
menerapkan arsitektur MVC memberikan pemahaman lebih mendalam mengenai
pentingnya struktur kode yang \textit{scalable} dan \textit{maintainable} dalam
proyek jangka panjang. Total 33 \textit{task} perbaikan dan peningkatan
\textit{widget} Flutter dikerjakan selama periode magang, mencakup penanganan
\textit{font size} kustom (\textit{TableFontUtils}), perbaikan
\textit{sorting}, \textit{null handling}, dan penyesuaian parameter SDUI dari
\textit{back-end}.

Mulai pertengahan Februari 2026, seiring transisi \textit{tech stack}
perusahaan, kemampuan membangun layar baru dari nol menggunakan React, Next.js,
TypeScript, Mantine, dan Tailwind berkembang cukup pesat dalam waktu yang
relatif singkat. Pengerjaan modul-modul dengan kompleksitas tinggi seperti Item
Cabang, Data Pasien, Dashboard Inventaris, dan komponen \textit{DataView}
(berbasis AG Grid) melatih kemampuan berpikir secara modular dan
\textit{component-driven}, serta membangun kepekaan terhadap konsistensi desain
di tingkat sistem.

Perkembangan paling signifikan terjadi mulai pertengahan April 2026, ketika
peran bergeser ke arah \textit{Full-Stack} melalui proses \textit{alignment}
layar-layar React dengan API dari \textit{back-end} Toraja (Django/Python).
Pada fase ini, pemahaman terhadap arsitektur sistem secara \textit{end-to-end}
-- mulai dari pembuatan dan modifikasi \textit{endpoint}, pengolahan
\textit{response}, navigasi \textit{iframe} antara React dan Toraja, hingga
distribusi data ke komponen \textit{front-end} -- berkembang jauh lebih cepat
dibandingkan pengerjaan \textit{front-end} pada umumnya. Modul Dashboard
Inventaris (Task \#13999) menjadi titik paling representatif dari pergeseran
peran ini, karena melibatkan pembuatan \textit{endpoint} baru secara aktif di
sisi Toraja.

Dari sisi non-teknis, kemampuan bekerja dalam tim lintas fungsi juga meningkat,
termasuk koordinasi aktif dengan Tim \textit{Back-End}, QA, dan \textit{Project
    Manager}. Penggunaan OpenProject untuk manajemen \textit{task} dan KPI serta
GitLab untuk \textit{version control} dan \textit{merge request} memperkuat
pemahaman terhadap alur kerja profesional dalam industri pengembangan perangkat
lunak, termasuk kedisiplinan dalam dokumentasi \textit{task} dan pencapaian
target poin bulanan.